\documentclass{article}
% Formal certificate:
% paper18_ufts008_final_conditional_certificate_closes_unified_field_theory_candidate_synthesis_theorem
\usepackage[margin=1in]{geometry}
\usepackage{array}
\usepackage{hyperref}
\usepackage{longtable}

\title{Unified Field Theory Candidate Synthesis From Finite-Capacity Causal Networks}
\author{Abraham Greenman}
\date{2026}

\newcommand{\finalcert}{\path{paper18_ufts008_final_conditional_certificate_closes_unified_field_theory_candidate_synthesis_theorem}}
\newcommand{\paperseventeen}{\path{paper17_ppa008_final_conditional_certificate_closes_physical_promotion_attempt_theorem}}

\begin{document}
\emergencystretch=3em
\maketitle

\begin{abstract}
This paper closes Paper 18 in the finite-capacity causal-network theorem
ladder.  Starting from the frozen Paper 17 conditional physical promotion
attempt interface theorem, it constructs a finite, network-native candidate
synthesis interface.  The result is conditional and deliberately
non-promoting: it closes a candidate-synthesis interface theorem, not a
completed unified field theory, not physical validation, not empirical
adequacy, not physical nature realization, and not physical promotion.
\end{abstract}

\section{Status and Scope}

The Paper 18 theorem is closed conditionally by the final certificate
\finalcert.  The closed object is an auditable finite interface for recording
candidate synthesis rows downstream of Paper 17.  The object is not a
completed physical theory and does not assert that nature is described by the
finite-capacity causal-network framework.

The claim boundary is part of the theorem.  The following claims remain false:
candidate synthesis success, completed unified field theory, physical nature
realization, physical promotion attempt success, physical promotion, physical
validation, empirical adequacy, certificate recovery, review acceptance,
reproduction success, protocol recovery, benchmark success, prediction success,
falsification success, observed particle catalog recovery, Standard Model
recovery, physical matter-field recovery, physical gauge-field recovery,
physical quantum dynamics, continuum quantum field theory, simulation-only
promotion, and fit-only calibration.

\section{Frozen Upstream Input}

Paper 18 consumes the closed Paper 17 conditional physical promotion attempt
interface theorem.  The frozen upstream references are:

\begin{itemize}
\item Paper 17 commit:
\texttt{5fa387b8aa40ffce1b930b68f6799c37083e21b2}.
\item Paper 17 formal endpoint:
\texttt{Paper17PhysicalPromotionAttemptTheoremContract.closed}.
\item Paper 17 final certificate: \paperseventeen.
\end{itemize}

The upstream import is intentionally narrow.  It imports the availability of a
closed conditional interface theorem through Paper 17; it does not import
physical promotion attempt success, physical promotion, physical validation,
empirical adequacy, physical nature realization, or a unified field theory.

\section{Theorem}

\textbf{Theorem 18.}  Given the frozen Paper 17 conditional physical promotion
attempt interface theorem and its upstream chain, there exists a finite,
bounded, auditable, network-native candidate-synthesis interface whose records
can name candidate rows, upstream dependencies, gate references, unresolved
risks, conflicts, audit events, rollback behavior, and the final conditional
certificate, while preserving the non-promotion and non-physical-success claim
boundary.

\textbf{Formal endpoint.}  The Lean theorem contract and closing certificate are:
\begin{quote}
\raggedright
\path{Paper18UnifiedFieldTheoryCandidateSynthesisTheoremContract.closed}

\finalcert
\end{quote}

\textbf{Non-goal.}  The theorem does not prove a completed unified field
theory.  It proves that the candidate-synthesis interface can be stated as
finite rows with fail-closed audit guards and no hidden promotion claim.

\section{Finite Interface Objects}

The construction uses only finite records.  In the formal scaffold these are
encoded as theorem contracts whose closure is a conjunction of finite-label,
bounded-row, auditability, and no-hidden-claim conditions.  In the Rust audit
surface, the same boundary is mirrored by bounded static record arrays and
explicit false success flags.

\subsection{Candidate Synthesis Records}

A candidate synthesis record contains a finite candidate identifier, a finite
synthesis-scope label, finite upstream dependency labels, finite gate-reference
labels, finite unresolved-risk labels, finite conflict labels, and a finite
audit-status descriptor.  Its audit status is
\texttt{bounded-auditable-nonpromoting}.  The record explicitly carries the
Paper 18 claim boundary, so the record cannot be interpreted as synthesis
success, a completed unified field theory, physical validation, empirical
adequacy, or promotion.

\subsection{Assumption, Dependency, and Gate References}

The second interface layer records finite assumption descriptors, dependency
descriptors, gate-reference descriptors, and upstream reference labels.  These
descriptors make the candidate row auditable, but they do not discharge the
assumptions, do not prove dependency sufficiency, and do not claim gate
success.

\subsection{Consistency, Conflict, and Residual Risk}

The third interface layer records finite consistency-check descriptors, finite
conflict rows, finite residual-risk rows, and explicit unresolved-status
labels.  The row says that the checks, conflicts, and risks are addressable as
finite audit objects.  It does not say that consistency has succeeded, that a
conflict is resolved, or that a residual risk is discharged.

\subsection{Paper 17 Compatibility}

The compatibility layer binds the Paper 18 interface to the exact frozen Paper
17 commit, endpoint, and final certificate.  Compatibility means that Paper 18
can reference the closed Paper 17 conditional interface.  It does not mean
physical promotion attempt success, physical promotion, physical validation,
empirical adequacy, physical nature realization, or a completed unified field
theory.

\subsection{Stability, Auditability, and Rollback}

The stability layer supplies finite candidate-revision labels, finite audit
event labels, finite rollback-trigger labels, and an explicit rollback behavior:
\texttt{bounded-revert-to-prior-audited-interface-row}.  Rollback is available
when the audit detects a missing upstream reference, unbounded row, hidden
promotion, or claim-boundary violation.  This is audit stability for interface
records only; it is not empirical stability, falsification success, or a
physical failure claim.

\subsection{No-Hidden-Claim Audit}

The no-hidden-claim audit covers docs, the Lean formal scaffold, Rust guards,
and upstream bindings.  It blocks hidden imports of unified-field success,
physical nature, physical validation, empirical adequacy, physical promotion
attempt success, physical promotion, simulation-only promotion, and fit-only
calibration.

\section{Rung Closure}

\begin{longtable}{>{\raggedright\arraybackslash}p{0.16\linewidth} >{\raggedright\arraybackslash}p{0.72\linewidth}}
\textbf{Rung} & \textbf{Closed content} \\
\hline
\texttt{UFTS-001} & Frozen upstream binding through Paper 17 and the claim-boundary scaffold. \\
\texttt{UFTS-002} & Finite candidate-synthesis records as bounded, auditable, non-promoting interface rows. \\
\texttt{UFTS-003} & Finite assumption, dependency, and gate-reference descriptors without assumption discharge, dependency sufficiency, or gate success. \\
\texttt{UFTS-004} & Finite consistency, conflict, residual-risk, and unresolved-status descriptors without consistency success, conflict-resolution success, or risk discharge. \\
\texttt{UFTS-005} & Paper 17 compatibility rows pinned to the frozen commit, endpoint, and certificate without importing promotion attempt success, promotion success, validation success, or empirical adequacy. \\
\texttt{UFTS-006} & Candidate revision, audit event, rollback trigger, and rollback behavior rows with rollback availability and without empirical stability, falsification success, or physical failure claims. \\
\texttt{UFTS-007} & No-hidden-claim audit over docs, formal code, Rust guards, and upstream bindings. \\
\texttt{UFTS-008} & Final conditional certificate closing the Paper 18 candidate-synthesis theorem contract. \\
\end{longtable}

\section{Proof}

The proof is by finite contract assembly.

\begin{enumerate}
\item \texttt{UFTS-001} establishes that the upstream chain through Paper 17 is
frozen and that no physical or unified-field success claim is imported.
\item \texttt{UFTS-002} constructs the candidate-synthesis record interface as
finite rows and proves that the rows are bounded, auditable, and
non-promoting.
\item \texttt{UFTS-003} adds finite assumption, dependency, and gate-reference
descriptors, while preserving the absence of assumption discharge, dependency
sufficiency, and gate success.
\item \texttt{UFTS-004} adds finite consistency, conflict, residual-risk, and
unresolved-status descriptors, while preserving the absence of consistency
success, conflict-resolution success, and risk discharge.
\item \texttt{UFTS-005} proves compatibility with the frozen Paper 17
conditional interface and blocks every physical promotion, validation, and
empirical-adequacy import.
\item \texttt{UFTS-006} adds finite audit and rollback behavior, preserving the
absence of empirical stability, falsification success, and physical failure
claims.
\item \texttt{UFTS-007} closes the no-hidden-claim audit across docs, Lean,
Rust, and upstream bindings.
\item \texttt{UFTS-008} assembles the final certificate as a conjunction of the
closed rungs \texttt{UFTS-001} through \texttt{UFTS-007}, the conditionality of
the certificate, and the explicit false physical-success and unified-field
success claims.
\end{enumerate}

The final Lean proof unfolds the Paper 18 theorem contract, unfolds the final
candidate-synthesis certificate, and simplifies the closed rung contracts.
Since each conjunct is supplied by a closed rung or by an explicit claim-boundary
guard, the theorem contract closes.

\section{Implementation and Reproducibility}

The formal artifact is:
\[
\texttt{GPD/formal/FiniteCapacity/UnifiedFieldTheoryCandidateSynthesis.lean}.
\]
The Rust audit surface is:
\[
\texttt{rust/cclab\_accel/src/lib.rs}
\]
with regression tests in:
\[
\texttt{rust/cclab\_accel/tests/unified\_field\_theory\_candidate\_synthesis.rs}.
\]

The verification commands are:

\begin{verbatim}
make test-fast
make lean-build
\end{verbatim}

The tests verify that each rung closes in order, that the final certificate
closes the Paper 18 theorem conditionally, that the repository state records
the final closure, and that the claim boundary remains non-promoting.

\section{Conclusion}

Paper 18 closes a finite candidate-synthesis interface theorem.  The final
certificate proves that candidate-synthesis records, dependency maps, gate
references, conflict rows, residual-risk rows, Paper 17 compatibility,
rollback behavior, and no-hidden-claim audits can be assembled as finite,
auditable, network-native rows.  The result is intentionally conditional and
non-promoting.  It does not assert a completed unified field theory or any
physical validation claim.

\end{document}
